\documentclass[11pt]{article}
\usepackage[margin=1in]{geometry}
\usepackage{amsmath,amssymb,amsthm,mathtools}
\usepackage[T1]{fontenc}
\usepackage{lmodern,microtype}
\usepackage[colorlinks=true,linkcolor=blue,urlcolor=blue,citecolor=blue]{hyperref}
\newtheorem{theorem}{Theorem}[section]
\newtheorem{lemma}[theorem]{Lemma}
\newtheorem{proposition}[theorem]{Proposition}
\newtheorem{remark}[theorem]{Remark}
\newcommand{\E}{\mathbb E}
\newcommand{\Pp}{\mathbb P}
\newcommand{\R}{\mathbb R}
\newcommand{\ind}{\mathbf1}
\newcommand{\F}{\mathcal F}
\newcommand{\G}{\mathcal G}
\newcommand{\N}{\mathcal N}
\newcommand{\code}[1]{\texttt{\small #1}}
\title{Brownian sticky clocks: construction, holding time\protect\\
and causal recovery in the observed filtration}
\author{Project Stardust\\\small A companion to a frozen verified library}
\date{October 29, 2026}
\begin{document}
\maketitle
\begin{abstract}
We reconstruct the complete construction of a Brownian sticky time
change from smooth convex approximations to absolute value. A compact
probability estimate and a summable subsequence give a measurable
continuous increasing mass, whose Stieltjes measure is supported on
Brownian zeros. Its inverse clock has an exact physical holding-time
identity. The clock is recoverable from the observed path up to the
current time, so both the process and its compensated square are
martingales in explicitly specified observed filtrations. At every
positive parameter and deterministic horizon, expected holding is
positive and the second moment is strictly smaller than elapsed time.
An independent Brownian source on the complementary clock yields two
genuine Brownian coordinate processes whose covariance equals expected
holding. Their finite-increment laws are proved by stopped Fourier
products, rather than inferred from Gaussian one-time marginals.
These classical mechanisms are presented with full proofs and exact
formal-source correspondences. The construction does not use the
separately unfinished closed-generator or stochastic-flow identification.
\end{abstract}

\section{The construction and its precise conclusions}
Let $B=(B_s)_{s\geq0}$ be standard real Brownian motion, $B_0=0$,
with a measurable realization in $C([0,\infty),\R)$. We use its raw
natural filtration, augmented by ambient measurable null events:
\[
 \F_s=\sigma(B_u:0\leq u\leq s)\vee\N.
\]
Here $\N$ consists of the null events in the ambient probability space;
this convention suffices for all modifications below. No right-continuous
enlargement of the filtration is implicit. The Brownian stochastic
integral is its ordinary predictable $L^2(ds\,d\Pp)$ isometry, with
conditional expectations at deterministic horizons. These standard
Brownian foundations are also constructed in the frozen source library.

\begin{theorem}\label{thm:main}
There is a measurable continuous nonnegative increasing process $L$,
starting from zero, with the following properties. For each fixed $s$,
\begin{equation}\label{eq:tanakaL2}
 L_s=|B_s|-\int_0^s\operatorname{sgn}(B_u)\,dB_u\quad\hbox{in }L^2.
\end{equation}
Its Stieltjes measure $dL$ is concentrated on $\{s:B_s=0\}$.
For every $\kappa\geq0$, set
\begin{equation}\label{eq:clocks}
 A_\kappa(s)=s+\kappa L_s,\qquad
 C_\kappa=A_\kappa^{-1},\qquad X_t^\kappa=B_{C_\kappa(t)}.
\end{equation}
The clock and process have measurable continuous path realizations.
For almost every sample, simultaneously for all $\kappa,t\geq0$,
\begin{equation}\label{eq:holding}
 H_\kappa(t):=\int_0^t\ind_{\{X_r^\kappa=0\}}\,dr
 =\kappa L_{C_\kappa(t)}=t-C_\kappa(t).
\end{equation}
For each fixed $\kappa$, write
\[
 \G_t^0=\sigma(X_r^\kappa:0\leq r\leq t),\qquad
 \G_t=\G_t^0\vee\N.
\]
The operational clock $C_\kappa(t)$ is $\G_t$-measurable. The process
$X^\kappa$ is a martingale in both $\G^0$ and $\G$, and
$(X_t^\kappa)^2-C_\kappa(t)$ is a martingale in $\G$.
For $\kappa,t>0$,
\begin{equation}\label{eq:variance}
 \E H_\kappa(t)>0,\qquad
 \E (X_t^\kappa)^2+\E H_\kappa(t)=t,\qquad
 \E (X_t^\kappa)^2<t.
\end{equation}
In particular, $\Pp(H_\kappa(t)>0)>0$.
\end{theorem}

The last assertion is positive probability, rather than an assertion
that holding is positive on every sample or that a terminal atom has
been computed at each time. At $\kappa=0$, $C_0(t)=t$ and $X^0=B$.
All clocks in the theorem use the same Brownian sample. Independence
of the clock and the Brownian motion is neither true nor needed.

The time-change construction is classical. The appendix of
Ayala--Carinci--Redig \cite{ACR}, equations (112)--(114), explicitly
uses Brownian time plus a local-time term and its inverse to construct
two-sided sticky Brownian motion. Bass \cite{Bass}, Section 2,
equations (2.3)--(2.4), places inverse additive clocks in the classical
speed-measure construction. Our contribution is an exposition of the
verified construction, its versions and observed-filtration statements;
we make no independent priority claim for these mechanisms.

\section{Smooth curvature masses and their fixed-time limit}
For $e>0$ let
\[
 r_e(x)=\sqrt{x^2+e^2},\qquad
 a_e(x)=\frac{x}{\sqrt{x^2+e^2}},\qquad
 b_e(x)=\frac{e^2}{(x^2+e^2)^{3/2}}.
\]
Direct differentiation gives $r_e'=a_e$, $r_e''=b_e$. Moreover,
\begin{equation}\label{eq:smoothbounds}
 0\leq r_e(x)-|x|\leq e,\quad |a_e(x)|\leq1,
 \quad 0\leq b_e(x)\leq e^{-1}.
\end{equation}
Indeed, $|x|\leq r_e(x)\leq|x|+e$ follows by squaring
nonnegative quantities. The slope bound follows from $|x|\leq r_e$;
the curvature bound follows from $r_e\geq e$.
Fix the positive sequence $e_n\downarrow0$ used in the source, and put
\begin{equation}\label{eq:approximants}
 L_t^n=\frac12\int_0^t b_{e_n}(B_u)\,du,
 \qquad M_t^n=r_{e_n}(B_t)-e_n-L_t^n.
\end{equation}
The first process is continuous, increasing and nonnegative on every
continuous Brownian path. The second is continuous and, by the smooth
It\^o formula,
\begin{equation}\label{eq:smoothito}
 M_t^n=\int_0^t a_{e_n}(B_u)\,dB_u
 \quad\hbox{at each fixed time, almost surely}.
\end{equation}
The coefficients in this formula are bounded and predictable. The
curvature term is integrable on every bounded horizon by
\eqref{eq:smoothbounds}; thus the formula involves no limiting local
time as an input.

Define $\operatorname{sgn}(0)=0$ and let
$M_t=\int_0^t\operatorname{sgn}(B_u)\,dB_u$, initially as an $L^2$
class for each $t$. The coefficient is an actual predictable field:
it is the bounded pointwise limit of the predictable continuous
functions $a_{e_n}(B_u)$. At every $x$,
$a_{e_n}(x)\to\operatorname{sgn}(x)$, and $|a_{e_n}|\leq1$.
Dominated convergence on $[0,T]\times\Omega$ and the stochastic
isometry give
\begin{equation}\label{eq:fieldlimit}
 \int_0^T\E|a_{e_n}(B_u)-\operatorname{sgn}(B_u)|^2du\to0,
 \qquad M_T^n\longrightarrow M_T\quad\hbox{in }L^2.
\end{equation}
Since $|r_{e_n}(B_T)-e_n-|B_T||\leq e_n$,
\begin{equation}\label{eq:fixedlimit}
 L_T^n\longrightarrow \ell_T:=|B_T|-M_T\quad\hbox{in }L^2.
\end{equation}
The closed positive cone in $L^2$ shows $\ell_T\geq0$ almost surely.
Similarly, $\ell_s\leq\ell_t$ as $L^2$ classes whenever $s\leq t$.
These fixed-time statements do not yet construct a continuous
increasing random path. That distinction motivates the next step.

\section{A continuous increasing mass from a summable subsequence}
We first give the maximal estimate needed to keep the limiting mass
on the original sample space.

\begin{lemma}\label{lem:maximal}
For $n,m$, $T\geq0$ and $\varepsilon>0$,
\begin{equation}\label{eq:maximal}
 \varepsilon\Pp\left(\sup_{u\leq T}|M_u^n-M_u^m|>
 \varepsilon\right)
 \leq \left(\int_0^T\E|a_{e_n}(B_u)-a_{e_m}(B_u)|^2du\right)^{1/2}.
\end{equation}
Consequently, $L^n$ is Cauchy in probability uniformly on each compact
time interval.
\end{lemma}
\begin{proof}
For a finite ordered set of times including $T$, $D=M^n-M^m$ is a
martingale. Partition the first-crossing event $\max|D|>
\varepsilon$ according to its first crossing time $u$. Each piece
belongs to $\F_u$, and
$|D_u|\leq\E(|D_T|\mid\F_u)$. Summing the integrals over the disjoint
pieces gives
$\varepsilon\Pp(\max|D|>\varepsilon)\leq\E|D_T|$.
An increasing collection of finite sets dense in $[0,T]$ gives the
continuous-time estimate by continuity and monotone convergence of
events. The stochastic isometry and Cauchy--Schwarz give
\eqref{eq:maximal}. This argument also proves measurability of the
crossing event, using only countably many evaluations.

From \eqref{eq:approximants} and the normalized absolute-value bound,
\begin{equation}\label{eq:massdifference}
 \sup_{u\leq T}|L_u^n-L_u^m|
 \leq e_n+e_m+\sup_{u\leq T}|M_u^n-M_u^m|.
\end{equation}
The right side tends to zero in probability by
\eqref{eq:fieldlimit} and \eqref{eq:maximal}. This proves the claim.
\end{proof}

For each $j\geq0$, choose $N_j$ such that for $n,m\geq N_j$,
\[
 \Pp\left(\sup_{u\leq j+1}|L_u^n-L_u^m|>2^{-j}\right)<2^{-j}.
\]
Set $n_0=N_0$ and $n_{j+1}=\max\{n_j+1,N_{j+1}\}$.
The consecutive bad events have summable probabilities, so
Borel--Cantelli implies that, almost surely, for all sufficiently
large $j$,
\begin{equation}\label{eq:fast}
 \sup_{u\leq j+1}|L_u^{n_{j+1}}-L_u^{n_j}|\leq2^{-j}.
\end{equation}
Fix any compact horizon $T$. For $k>j$ sufficiently large that
$j+1\geq T$, telescoping bounds the uniform difference by
$\sum_{i=j}^{k-1}2^{-i}$. Completeness of $C([0,T])$ gives a uniform
limit. These limits agree on overlapping compact intervals, giving
one continuous path $L$ on $[0,\infty)$. It is increasing,
nonnegative and zero at the origin, because every approximant has
these properties. The same probability-one event works for all
compact horizons, by taking integer horizons first.

The event in \eqref{eq:fast} is measurable, since uniform crossings
of continuous paths are countable evaluation events. Each map
$\omega\mapsto L^{n_j}(\omega)$ is measurable into the continuous
path space with its compact-convergence topology: evaluations are
measurable, and on this separable path space they generate its Borel
sigma field. Set the limiting path equal to zero on the exceptional
event. The resulting random continuous path is measurable, since it
is the path-space limit on a measurable event of convergence. This
normalization gives a continuous increasing mass starting at zero on
every sample. The source instead centers its Stieltjes cumulative
function on exceptional samples; both conventions agree almost surely.

At each fixed $T$, almost-sure subsequence convergence to $L_T$
and \eqref{eq:fixedlimit} identify
\begin{equation}\label{eq:coordinate}
 L_T=\ell_T\quad\hbox{almost surely},\qquad L_T\in L^2.
\end{equation}
For example, take a further subsequence along which the $L^2$
convergence is almost sure, and use uniqueness of real limits.
Thus \eqref{eq:tanakaL2} holds. Define the continuous process
$\widehat M_t=|B_t|-L_t$. At each fixed time it is a version of
$M_t$; conditional-expectation identities transfer to this version.
It is adapted to $\F$ and a martingale there. The identity
\begin{equation}\label{eq:alltanaka}
 |B_t|=\widehat M_t+L_t
\end{equation}
is now a simultaneous pathwise identity, rather than a substitution
of a random time into unrelated fixed-time null-set statements.

\section{Support of the Stieltjes mass}
For a compact interval $[s,t]$ on which $|B_u|\geq\delta>0$,
\[
 0\leq L_t^n-L_s^n
 =\frac12\int_s^t\frac{e_n^2}{(B_u^2+e_n^2)^{3/2}}du
 \leq\frac{(t-s)e_n^2}{2\delta^3}.
\]
Uniform subsequence convergence therefore gives $L_t=L_s$.
Every compact interval free of Brownian zeros has a positive such
$\delta$, by continuity and compactness. Consequently the Stieltjes
measure $dL$ vanishes locally on the open set $\{B\ne0\}$.
Second countability supplies a countable cover by such neighborhoods,
so
\begin{equation}\label{eq:support}
 dL(\{s:B_s\ne0\})=0.
\end{equation}
Continuity of $L$ implies that $dL$ has no atoms, including at zero;
its interval masses are
$dL((s,t])=L_t-L_s$, and $dL([0,t])=L_t$.

Ordinary Brownian time gives no mass to the same zero set. Indeed,
$\Pp(B_s=0)=0$ for every $s>0$, since $B_s$ has a nondegenerate
Gaussian law. Tonelli yields
\[
 \E\int_0^N\ind_{\{B_s=0\}}ds
 =\int_0^N\Pp(B_s=0)ds=0.
\]
The integrand is nonnegative. Taking all integer $N$ gives,
on one probability-one event,
\begin{equation}\label{eq:brownianzero}
 ds(\{s:B_s=0\})=0.
\end{equation}
The two different measures in \eqref{eq:support} and
\eqref{eq:brownianzero} are exactly what creates physical holding.

Nontriviality follows without a distributional formula for local
time. The stochastic integral is centered, so
\begin{equation}\label{eq:expectedmass}
 \E L_T=\E|B_T|>0\qquad(T>0).
\end{equation}
Both terms are integrable by \eqref{eq:coordinate}. Positivity uses
only the nondegenerate Brownian marginal, which makes $|B_T|>0$
almost surely. We use no assertion that $L_T$ is positive on every
sample.

\section{Invertible clocks and exact measure transport}
We record the deterministic clock mechanism in a form that can be
checked before inserting any randomness.

\begin{lemma}\label{lem:inverse}
If $l:[0,\infty)\to[0,\infty)$ is continuous, increasing and
$l(0)=0$, then $A(s)=s+\kappa l(s)$ is a continuous strictly
increasing bijection of $[0,\infty)$ onto itself. Its inverse $C$
is continuous, increasing, and satisfies
\begin{equation}\label{eq:inverseproperties}
 C(t)\leq t,\qquad
 0\leq C(t)-C(r)\leq t-r\quad(r\leq t),\qquad
 C(t)+\kappa l(C(t))=t.
\end{equation}
\end{lemma}
\begin{proof}
For $v>u$, $A(v)-A(u)\geq v-u>0$. Also $A(0)=0$ and
$A(s)\geq s\to\infty$; the intermediate value theorem gives
surjectivity. The inverse of a continuous increasing bijection of an
interval is continuous and increasing. Substitute $C(t)$ into $A$
to get the balance identity and $C(t)\leq t$. Subtract the balance
identities at $r,t$; the nonnegative mass increment gives the upper
increment bound. The lower bound follows from monotonicity.
\end{proof}

Apply this to each mass sample. For fixed $\kappa,t$, measurability
of the inverse follows from
\begin{equation}\label{eq:inverseevent}
 \{C_\kappa(t)<b\}=\{t<A_\kappa(b)\},\qquad b\geq0.
\end{equation}
Each right side is measurable. Continuous time paths and measurable
coordinates give jointly measurable clocks and measurable continuous
path maps. Evaluation $(p,s)\mapsto p(s)$ is continuous for compact
convergence on bounded time neighborhoods. Hence
$X_t^\kappa=B_{C_\kappa(t)}$ has measurable coordinates, jointly
measurable time/sample evaluation, and a measurable continuous path
realization. No random time is selected using an independent copy of $B$.

On $[0,\infty)$ let $ds$ denote Lebesgue measure with no atom at
zero. For any continuous mass path from the construction, define
$\eta_\kappa=ds+\kappa\,dL$. Then
\begin{equation}\label{eq:pushforward}
 (A_\kappa)_\#\eta_\kappa=dt.
\end{equation}
To prove this, test $[0,t]$. Strict invertibility gives
$A_\kappa^{-1}([0,t])=[0,C_\kappa(t)]$, whose $\eta_\kappa$-mass
is $C_\kappa(t)+\kappa L_{C_\kappa(t)}=t$. These intervals form a
generating pi system and exhaust the half-line by finite-mass
integer intervals. Uniqueness of locally finite measures proves
\eqref{eq:pushforward} on every Borel set. This is a measure identity,
so it also applies to each sample's own Borel zero set.

\section{Holding time and a strict second-moment budget}
Fix a sample on which \eqref{eq:support} and
\eqref{eq:brownianzero} hold. For any $\kappa,t\geq0$, let
$Z=\{s:B_s=0\}$ and
$E_t=\{r\leq t:X_r^\kappa=0\}$. These sets are Borel, and
\[
 A_\kappa^{-1}(E_t)=Z\cap[0,C_\kappa(t)].
\]
By \eqref{eq:pushforward},
\begin{align*}
 \int_0^t\ind_{\{X_r^\kappa=0\}}dr
 &=\eta_\kappa(Z\cap[0,C_\kappa(t)])\\
 &=\kappa\,dL([0,C_\kappa(t)])
 =\kappa L_{C_\kappa(t)}=t-C_\kappa(t).
\end{align*}
The ordinary-time contribution is zero by
\eqref{eq:brownianzero}; concentration of $dL$ on $Z$ gives the
second line. The proof uses just two fixed probability-one path
properties. Thus it holds simultaneously for all uncountably many
$\kappa,t$ without an uncountable intersection of null sets.

For $\kappa>0$ and each sample, strict inversion also gives
\begin{equation}\label{eq:holdingpositive}
 t-C_\kappa(t)>0
 \ \Longleftrightarrow\ C_\kappa(t)<t
 \ \Longleftrightarrow\ t<A_\kappa(t)
 \ \Longleftrightarrow\ L_t>0.
\end{equation}
For $t>0$, \eqref{eq:expectedmass} implies
$\Pp(L_t>0)>0$. Since $0\leq H_\kappa(t)\leq t$, it is integrable;
\eqref{eq:holdingpositive} implies $\E H_\kappa(t)>0$.
The second-moment part of \eqref{eq:variance} will follow from
bounded sampling below. This separates its probabilistic input from
the pathwise measure calculation.

\section{Bounded sampling with the clock's own Brownian noise}
The continuous mass $L_s$ agrees at each fixed time with the raw
Brownian-past $L^2$ class in \eqref{eq:tanakaL2}. Hence $L_s$ is
$\F_s$-measurable under our ambient null augmentation. The same is
true of $A_\kappa(s)$. The inverse is a bounded stopping time because
\begin{equation}\label{eq:stopping}
 \{C_\kappa(t)\leq s\}=\{t\leq A_\kappa(s)\}\in\F_s,
 \qquad C_\kappa(t)\leq t.
\end{equation}
The weak inequality here, as opposed to the strict inequality in
\eqref{eq:inverseevent}, is required for the stopping-time sigma field.

We give the bounded sampling argument used in the construction.
If $Y$ is a continuous martingale and $\sigma\leq\tau\leq T$
are stopping times, with integrable compact-time maximum, then
\begin{equation}\label{eq:optionalsampling}
 \E(Y_\tau\mid\F_\sigma)=Y_\sigma.
\end{equation}
Round both times up to a dyadic grid, obtaining finite-valued
stopping times $\sigma_n\leq\tau_n\leq T+1$. For a finite grid,
\eqref{eq:optionalsampling} follows by writing the sampled difference
as a sum of martingale increments multiplied by past-measurable
indicators. For any $E\in\F_\sigma$ we also have
$E\in\F_{\sigma_n}$, so
$\E\ind_EY_{\tau_n}=\E\ind_EY_{\sigma_n}$.
Path continuity gives almost-sure convergence of both values.
Their absolute values are bounded by
$\sup_{u\leq T+1}|Y_u|$, so dominated convergence passes this equality
to the limit. Stopped values are $\F_\sigma$-measurable: a continuous
adapted process is progressively measurable, and its evaluation at
a stopping time is measurable in the stopping-time sigma field.
This last assertion follows by approximating the continuous path on
each deterministic interval by time-step functions, then restricting
to $\{\sigma\leq s\}$. Integrability and the integral identity for
every $E$ prove \eqref{eq:optionalsampling}.

The maximum condition holds for the processes used here. The $L^2$
martingales $B$ and $\widehat M$ have integrable compact maxima by
the finite-grid $L^2$ maximal inequality and path continuity.
For completeness, that inequality follows by the same first-crossing
partition with levels $a>0$:
$a\Pp(Y^*>a)\leq\E(|Y_T|\ind_{Y^*>a})$.
Integrate twice the resulting inequality in $a$, then use
Cauchy--Schwarz to get
$\E(Y^*)^2\leq2\E(|Y_T|Y^*)\leq
2\|Y_T\|_2\|Y^*\|_2$, first with truncated maxima and then by
monotone convergence. Thus $\|Y^*\|_2\leq2\|Y_T\|_2$.
Finally,
\[
 \sup_{u\leq T}|B_u^2-u|\leq\sup_{u\leq T}B_u^2+T
\]
is integrable. So sampling applies also to the martingale $B_s^2-s$.

Let $\F_{C_\kappa(t)}$ be the stopping-time sigma field and denote
this increasing physical filtration by $\mathcal H_t$. Applying
\eqref{eq:optionalsampling} at $C_\kappa(r)\leq C_\kappa(t)$ gives
martingales
\begin{equation}\label{eq:physicalmart}
 X_t^\kappa,\qquad
 (X_t^\kappa)^2-C_\kappa(t),\qquad
 \widehat M_{C_\kappa(t)}
 \quad\hbox{in }\mathcal H.
\end{equation}
Their sampled values are adapted by stopped-value measurability.
Sampling the centered square martingale gives
$\E(X_t^\kappa)^2=\E C_\kappa(t)$; its integrability follows from
the preceding compact-maximum bound. Combine this equality with
\eqref{eq:holding} to obtain \eqref{eq:variance}, including its strict
inequality. The process is centered by the same argument applied to
$B$, so its second moment is indeed its variance.

The pathwise Tanaka identity also becomes
\[
 |X_t^\kappa|=\widehat M_{C_\kappa(t)}+L_{C_\kappa(t)},
 \qquad \kappa L_{C_\kappa(t)}=t-C_\kappa(t).
\]
This uses \eqref{eq:alltanaka} on the common path event. Substitution
of the random time would not be justified from fixed-time
$L^2$ identities alone.

\section{Causal clock recovery and the observed martingales}
For a continuous real path $p$ define
\begin{equation}\label{eq:pathfunctional}
 h(p,t)=\int_0^t\ind_{\{p(r)=0\}}dr,\qquad
 c(p,t)=t-h(p,t).
\end{equation}
Then $0\leq h(p,t)\leq t$ and $0\leq c(p,t)\leq t$ for every
path. Both functionals are jointly Borel measurable in $(p,t)$.
Indeed,
\[
 (p,t,r)\longmapsto
 \ind_{\{0<r\leq t\}}\ind_{\{p(r)=0\}}
\]
is Borel measurable, since evaluation is continuous and the indicated
sets are Borel. Integrating this nonnegative function in $r$ proves
measurability by the measurable-integral theorem. The integral is
finite, so subtraction in \eqref{eq:pathfunctional} is legitimate.
Both functionals depend only on $p|_{[0,t]}$.

To make the raw-past measurability explicit, the map
$\omega\mapsto X^\kappa(\omega)|_{[0,t]}$ into $C([0,t])$ is
$\G_t^0$-measurable: its evaluations are precisely the coordinates
generating $\G_t^0$, and they generate the Borel sigma field of
$C([0,t])$. Extend a prefix constantly after $t$; this is a measurable
map back into the full continuous path space. Applying
\eqref{eq:pathfunctional} to the extension therefore gives a
$\G_t^0$-measurable version of the recovered time. By
\eqref{eq:holding},
\begin{equation}\label{eq:recovery}
 c(X^\kappa,t)=C_\kappa(t)\quad\hbox{almost surely}.
\end{equation}
The actual clock coordinate is ambient measurable. A measurable
variable equal almost surely to a raw-past measurable one is
measurable in $\G_t^0\vee\N$. Thus the actual clock is adapted to
$\G$, with a causal raw-past version. We do not silently identify
these two assertions on unaugmented sigma fields.

For $u\leq t$, the coordinate $X_u^\kappa$ belongs to
$\mathcal H_u\subseteq\mathcal H_t$, so
$\G_t^0\subseteq\mathcal H_t$. The elementary shrinking-filtration
argument now transfers the first martingale in
\eqref{eq:physicalmart}: for $r\leq t$,
\[
 \E(X_t^\kappa\mid\G_r^0)
 =\E\bigl(\E(X_t^\kappa\mid\mathcal H_r)\mid\G_r^0\bigr)
 =\E(X_r^\kappa\mid\G_r^0)=X_r^\kappa.
\]
It remains a martingale after ambient null augmentation, because
adjoining ambient null events changes neither integrals nor
conditional expectations up to almost-sure equality.

For the compensated square, use the augmented physical filtration
$\mathcal H_t\vee\N$. Then $\G_t\subseteq\mathcal H_t\vee\N$,
and $(X_t^\kappa)^2-C_\kappa(t)$ is $\G_t$-measurable by
\eqref{eq:recovery}. The same tower argument transfers the second
martingale in \eqref{eq:physicalmart} to $\G$. This completes every
assertion of Theorem~\ref{thm:main}.

\section{A balanced coupling of two genuine Brownian coordinates}
The same inverse clock has a useful consequence beyond the single
sticky path. On the product of two independent Brownian source spaces,
let $W$ be independent of the Brownian motion $B$ used to construct
$L$ and $C_\kappa$. For the rest of this section fix $\kappa$ and abbreviate
\[
 C(t)=C_\kappa(t),\qquad D(t)=2t-C(t),\qquad
 S_t=B_{C(t)},\quad V_t=W_{D(t)}.
\]
The inverse increment bound in \eqref{eq:inverseproperties} implies
that $D$ is continuous and increasing, with $t\leq D(t)\leq2t$,
$D(0)=0$, and $C(t)+D(t)=2t$. Define
\begin{equation}\label{eq:pair}
 Y_t^+=\frac{V_t+S_t}{\sqrt2},\qquad
 Y_t^-=\frac{V_t-S_t}{\sqrt2}.
\end{equation}
These are measurable continuous paths, by measurable evaluation and
the product probability law. The source centers both Brownian paths
at their initial values on exceptional samples; they start at zero
almost surely, so this convention gives the same path law.

\begin{theorem}\label{thm:pair}
Each coordinate $Y^+$ and $Y^-$ has the law of standard Brownian
motion, including all finite collections of independent Gaussian
increments. For every $t\geq0$,
\begin{equation}\label{eq:paircov}
 \operatorname{Cov}(Y_t^+,Y_t^-)
 =\E H_\kappa(t)
 =\E\int_0^t\ind_{\{Y_r^+=Y_r^-\}}dr.
\end{equation}
For $\kappa,t>0$ the covariance is strictly positive, so the terminal
coordinates are not independent. For every $\kappa,T>0$ there exists
a deterministic $s\in(0,T]$ for which
\begin{equation}\label{eq:diagonal}
 \Pp(Y_s^+=Y_s^-)>0.
\end{equation}
At such $s$, the joint law differs from the product of its Gaussian
marginals, which gives zero mass to the diagonal.
\end{theorem}

The assertion that each coordinate is Brownian concerns its own full
path law. We do not deduce a Brownian property relative to the joint
pair filtration merely from the marginal assertion. Nor do we
assert a stochastic flow from this single two-point construction.

\subsection*{Why a single Gaussian marginal does not prove the theorem}
Given the entire $B$ sample, the clock $D(t)$ is deterministic and
$W_{D(t)}$ is conditionally Gaussian. But $B_{C(t)}$ is sampled at a
stopping time depending on $B$ itself. Treating it as an independent
Gaussian with variance $C(t)$ is invalid. The compensating Fourier
martingale, rather than that invalid conditional independence,
gives the cancellation in the proof below.

For $v\in\R$ define the complex process
\begin{equation}\label{eq:phase}
 Z_v(s)=\exp\bigl(v^2s/2+ivB_s\bigr).
\end{equation}
This is a continuous complex martingale on every bounded time
interval. Indeed, Brownian independent increments give
\[
 \E(e^{iv(B_t-B_s)}\mid\F_s)=e^{-v^2(t-s)/2},
\]
which verifies the martingale identity directly. Its modulus is the
deterministic quantity $e^{v^2s/2}$; its real and imaginary parts
therefore satisfy bounded sampling as in Section 7. It is never
zero and $|Z_v(s)|\geq1$.

\begin{lemma}\label{lem:products}
For ordered bounded Brownian stopping times
$\tau_0\leq\cdots\leq\tau_n\leq T$ and real $v_0,\ldots,v_{n-1}$,
\begin{equation}\label{eq:productphase}
 \E\prod_{j=0}^{n-1}\frac{Z_{v_j}(\tau_{j+1})}
                                  {Z_{v_j}(\tau_j)}=1.
\end{equation}
\end{lemma}
\begin{proof}
Write the product through $j=k-1$ as $R_k$, with $R_0=1$.
It is $\F_{\tau_k}$-measurable. Every factor has modulus
\[
 \exp\bigl(v_j^2(\tau_{j+1}-\tau_j)/2\bigr)
 \leq e^{v_j^2T/2},
\]
so the finite product is bounded by a deterministic constant. Also
$1/Z_{v_k}(\tau_k)$ is bounded and
$\F_{\tau_k}$-measurable. Bounded sampling of $Z_{v_k}$, separately
in real and imaginary parts, gives
\begin{align*}
 \E R_{k+1}
 &=\E\left[\frac{R_k}{Z_{v_k}(\tau_k)}
       \E(Z_{v_k}(\tau_{k+1})\mid\F_{\tau_k})\right]\\
 &=\E R_k.
\end{align*}
All products are integrable by the deterministic bounds. Induction
proves \eqref{eq:productphase}; independence of the stopping intervals
has not been used.
\end{proof}

\section{The full finite-increment law by exact compensation}
More generally let $a,b\in\R$ satisfy $a^2=b^2$, and put
$Q_t=aV_t+bS_t$. For deterministic ordered times
$0\leq t_0\leq\cdots\leq t_n$ and real $u_0,\ldots,u_{n-1}$,
we prove
\begin{equation}\label{eq:balancedincrements}
 \E\exp\left(i\sum_{j=0}^{n-1}u_j(Q_{t_{j+1}}-Q_{t_j})\right)
 =\prod_{j=0}^{n-1}e^{-(u_ja)^2(t_{j+1}-t_j)}.
\end{equation}
Condition first on the full $B$ sample. The $D(t_j)$ are then
ordered deterministic times for the independent Brownian motion $W$.
Its Gaussian independent increments give the conditional expectation
\begin{equation}\label{eq:pairfiber}
 \prod_{j=0}^{n-1}
 e^{-(u_ja)^2(D(t_{j+1})-D(t_j))/2}
 e^{iu_jb(B_{C(t_{j+1})}-B_{C(t_j)})}.
\end{equation}
This use of conditioning can equally be written as an iterated
integral on the product of the two source spaces. The original
integrand has modulus one, so Fubini is valid.

Write $\tau_j=C(t_j)$, $\Delta t_j=t_{j+1}-t_j$ and
$\Delta\tau_j=\tau_{j+1}-\tau_j$. Since
$\Delta D_j=2\Delta t_j-\Delta\tau_j$ and
$(u_ja)^2=(u_jb)^2$, each factor in \eqref{eq:pairfiber} is exactly
\begin{equation}\label{eq:exactcompensation}
 e^{-(u_ja)^2\Delta t_j}
 \exp\left((u_jb)^2\Delta\tau_j/2+
             iu_jb(B_{\tau_{j+1}}-B_{\tau_j})\right)
 =e^{-(u_ja)^2\Delta t_j}
  \frac{Z_{u_jb}(\tau_{j+1})}{Z_{u_jb}(\tau_j)}.
\end{equation}
The times $\tau_j$ are genuine ordered stopping times, all bounded
by $t_n$. Lemma~\ref{lem:products} shows that the expectation of
the product of the phase ratios is one. Pulling out the deterministic
factors proves \eqref{eq:balancedincrements}.

For $a=1/\sqrt2$ and $b=\pm1/\sqrt2$, the right side becomes
\[
 \prod_{j=0}^{n-1}\exp\bigl(-u_j^2\Delta t_j/2\bigr),
\]
the characteristic function of a vector of independent centered
Gaussian increments with variances $\Delta t_j$. Equality of
characteristic functions determines finite-dimensional probability
laws. Thus the increment vectors of $Y^+$ and $Y^-$ separately
have exactly those Brownian laws, including zero-length intervals.
Both paths are continuous and start at zero. This proves the full
Brownian-coordinate assertion in Theorem~\ref{thm:pair}, without
inferring it from a one-time Gaussian calculation.

The proof remains valid when $t_0>0$: the first phase ratio is
conditioned at $C(t_0)$, and the empty product at that time is one.
All parameter-dependent bounds are finite for the fixed collection
of frequencies and times. No limiting product, random signed law
or independence of the Brownian stopping-clock increments is used.

\section{Covariance, coincidence and the deterministic-time diagonal}
The elementary pathwise identity
\begin{equation}\label{eq:pairdifference}
 Y_t^+-Y_t^-=\sqrt2 S_t
\end{equation}
implies
\[
 \{Y_t^+=Y_t^-\}=\{S_t=0\},\qquad
 \int_0^t\ind_{\{Y_r^+=Y_r^-\}}dr=H_\kappa(t)
\]
on the same common event as \eqref{eq:holding}. The product-space
extension preserves that event because the second source is an
independent probability factor. Each Brownian coordinate is
centered with second moment $t$. From
\eqref{eq:pairdifference} and the exact single-process variance,
\[
 \E(Y_t^+-Y_t^-)^2=2\E S_t^2=2\E C(t).
\]
Expand the square. The product $Y_t^+Y_t^-$ is integrable by
Cauchy--Schwarz, so the calculation is valid and gives
\[
 2t-2\E(Y_t^+Y_t^-)=2\E C(t),\qquad
 \operatorname{Cov}(Y_t^+,Y_t^-)=t-\E C(t)=\E H_\kappa(t).
\]
This proves \eqref{eq:paircov} and, for $\kappa,t>0$, strict
positivity. Independence would force zero covariance for these
$L^2$ variables, so it is impossible at every such terminal time.

To prove the stronger diagonal statement at some deterministic
time in each horizon, Tonelli gives
\begin{equation}\label{eq:tonelliatom}
 \int_0^T\Pp(S_s=0)ds=\E H_\kappa(T)>0.
\end{equation}
Joint measurability of $S$ was proved in Section 5. If every
deterministic $s\in(0,T]$ had zero boundary probability, the integral
would vanish. Hence some such $s$ satisfies
$\Pp(S_s=0)>0$, and \eqref{eq:pairdifference} gives
\eqref{eq:diagonal}. The product of two nondegenerate Gaussian
marginals has zero diagonal mass: integrate the singleton mass
$\Pp(Y_s^-=x)=0$ over the first Gaussian marginal. Thus the
joint law cannot be their product. This completes
Theorem~\ref{thm:pair}.

The covariance conclusion holds at every positive deterministic
time, whereas the diagonal conclusion just proved is existential
within each positive horizon. The two quantifiers should not be
interchanged. Classical sticky Brownian couplings already combine
Brownian coordinates with coincidence occupation. Howitt--Warren
\cite{HW}, equations (1)--(2), formulate such couplings using a
common-filtration bracket and local time. The present source-mapped
claims establish the stated coordinate path laws, covariance and
coincidence facts; no identification of a joint bracket or of an
entire stochastic flow is silently added.

\section{What is established and what remains separate}
Clock recovery uses the observed prefix, not a terminal coordinate.
It does not construct a strong solution with a prescribed physical-time
Brownian driver. Bass's sticky SDE \cite{Bass} distinguishes weak
existence and uniqueness from strong existence; the operational
representation \eqref{eq:clocks} does not identify that driver filtration.

The source library separately proves a densely defined closable
continuous-function generator relation and a maximum-principle
bound. Its range condition and closed-semigroup identification are
not supplied by those results. No such missing step is assumed in
this paper. In particular, positive two-time conditional operators
alone do not establish Chapman--Kolmogorov consistency or an
all-start transition family. No stochastic flow, generator-domain
classification or uniqueness of the sticky martingale problem is
asserted by Theorem~\ref{thm:main}.

\section{Formal provenance and reproducibility}
All mathematical sources used here are frozen at
\begin{center}\code{bab688bb6140cd6316c80680095a914ea0d346db}.\end{center}
The theorem map records exact declarations, complete headers, source
hashes and baseline links for the mass, clock, martingale and pair
constructions. The original full-build and transitive-axiom receipt
and its referenced evidence hashes were checked. No mathematical
source was changed and no competing research build was run.

The printed proofs reorganize verified arguments rather than adding
formal declarations. Their distinctions between fixed-time $L^2$
classes, common continuous paths, raw-past versions and augmented
actual coordinates are retained. The pair proof supplies the full
finite-increment argument beyond Gaussian marginals. The cited
primary works provide classical antecedents; no new sticky process,
local-time mechanism or priority is claimed.

\begin{thebibliography}{9}
\bibitem{ACR} M. Ayala, G. Carinci and F. Redig,
\emph{Condensation of SIP particles and sticky Brownian motion},
Journal of Statistical Physics \textbf{183} (2021), article 40.
\href{https://doi.org/10.1007/s10955-021-02775-5}
{doi:10.1007/s10955-021-02775-5}.
\bibitem{Bass} R. F. Bass,
\emph{A stochastic differential equation with a sticky point},
Electronic Journal of Probability \textbf{19} (2014), paper 32, 1--22.
\href{https://doi.org/10.1214/EJP.v19-2350}{doi:10.1214/EJP.v19-2350}.
\bibitem{HW} C. Howitt and J. Warren,
\emph{Dynamics for the Brownian web and the erosion flow},
primary preprint (2007).
\href{https://arxiv.org/abs/math/0702542}{arXiv:math/0702542}.
\end{thebibliography}
\end{document}
